\chapter{Le seste eccedenti}
\label{chap:seste_eccedenti}

Le \textbf{seste eccedenti} sono accordi cromatici di preparazione alla dominante, più intensi della sesta napoletana perché la tensione è doppia: il basso scende di semitono e la voce superiore sale di semitono, convergendo entrambi sulla nota della dominante.

\section{Struttura}
\label{sec:struttura_seste}

\begin{definizione}[Sesta eccedente]
	\label{def:sesta_eccedente_cap}
	La \textbf{sesta eccedente} è un accordo cromatico che contiene l'intervallo di sesta eccedente tra il \textbf{VI grado abbassato} (al basso) e il \textbf{IV grado alzato}. In La minore: Fa--Re$\sharp$. La sesta Fa--Re$\sharp$ misura 10 semitoni, come una settima minore, ma è scritta come sesta: è questo che le dà il significato di accordo che si espande.
\end{definizione}

Per \textbf{moto contrario} i due suoni si allargano fino all'\textbf{ottava sul V grado}: il basso scende di semitono (Fa $\searrow$ Mi), la voce superiore sale di semitono (Re$\sharp$ $\nearrow$ Mi). L'accordo ha funzione di \textbf{predominante} e risolve sul V, direttamente o passando per il 6/4 di cadenza.

\begin{definizione}[Le tre specie]
	\label{def:tre_seste}
	In La minore, con il Fa al basso:
	\begin{itemize}
		\item \textbf{Italiana} (sesta it.): VI$\flat$--I--IV$\sharp$ $=$ Fa--La--Re$\sharp$. Tre suoni: si raddoppia la tonica La. Cifratura $\sharp$6.
		\item \textbf{Francese} (sesta fr.): VI$\flat$--I--II--IV$\sharp$ $=$ Fa--La--Si--Re$\sharp$. Cifratura $\sharp$6 4 3.
		\item \textbf{Tedesca} (sesta ted.): VI$\flat$--I--III$\flat$--IV$\sharp$ $=$ Fa--La--Do--Re$\sharp$. Cifratura $\sharp$6 5 3.
	\end{itemize}
	In tonalità maggiore si usano gli stessi accordi con il VI grado abbassato (Do maggiore: La$\flat$-Do-(Re oppure Mi$\flat$)-Fa$\sharp$).
\end{definizione}

\begin{esempio}[Le tre seste eccedenti in La minore]
	\label{ex:seste_costruzione_la}
	\musica{es_cap22_01}
\end{esempio}

\begin{esempio}[Le tre seste eccedenti in Do maggiore (VI grado abbassato: La$\flat$)]
	\label{ex:seste_costruzione_do}
	\musica{es_cap22_02}
\end{esempio}

Le note delle tre specie in tutte le tonalità (valgono sia per la tonalità maggiore sia per la minore omonima):

\begin{center}
\small
\begin{tabular}{|l|l|l|l|l|l|}
\hline
Tonalità & VI$\flat$ (basso) & I & III$\flat$ (solo ted.) & II (solo fr.) & IV$\sharp$ \\
\hline
Do & La$\flat$ & Do & Mi$\flat$ & Re & Fa$\sharp$ \\
Sol & Mi$\flat$ & Sol & Si$\flat$ & La & Do$\sharp$ \\
Re & Si$\flat$ & Re & Fa & Mi & Sol$\sharp$ \\
La & Fa & La & Do & Si & Re$\sharp$ \\
Mi & Do & Mi & Sol & Fa$\sharp$ & La$\sharp$ \\
Si & Sol & Si & Re & Do$\sharp$ & Mi$\sharp$ \\
Fa$\sharp$ & Re & Fa$\sharp$ & La & Sol$\sharp$ & Si$\sharp$ \\
Re$\flat$ & Si$\flat\flat$ & Re$\flat$ & Fa$\flat$ & Mi$\flat$ & Sol \\
La$\flat$ & Fa$\flat$ & La$\flat$ & Do$\flat$ & Si$\flat$ & Re \\
Mi$\flat$ & Do$\flat$ & Mi$\flat$ & Sol$\flat$ & Fa & La \\
Si$\flat$ & Sol$\flat$ & Si$\flat$ & Re$\flat$ & Do & Mi \\
Fa & Re$\flat$ & Fa & La$\flat$ & Sol & Si \\
\hline
\end{tabular}
\end{center}

\section{Risoluzione}
\label{sec:risoluzione_seste}

\begin{regola}[Condotta delle seste eccedenti]
	\label{reg:condotta_seste}
	\begin{itemize}
		\item Il basso (VI$\flat$) scende di semitono alla dominante; il IV$\sharp$ sale di semitono alla dominante.
		\item Non si raddoppiano né il VI$\flat$ né il IV$\sharp$: sono le note caratteristiche e devono risolvere in semitono. Nell'italiana si raddoppia la tonica.
		\item \textbf{Italiana}: la tonica raddoppiata si divide (La $\searrow$ Sol$\sharp$ e La $\nearrow$ Si). \textbf{Francese}: La $\searrow$ Sol$\sharp$, il Si resta fermo.
		\item \textbf{Tedesca}: se risolve direttamente sul V produce \textbf{quinte parallele} (Fa--Do $\rightarrow$ Mi--Si). Si passa quindi per il \textbf{6/4 di cadenza}: Do e La restano fermi, Re$\sharp$ e Fa si portano a Mi; poi La $\searrow$ Sol$\sharp$, Do $\searrow$ Si.
	\end{itemize}
\end{regola}

\begin{esempio}[Seste italiana e francese in La minore]
	\label{ex:italiana_francese_minore}
	Il Fa scende a Mi, il Re$\sharp$ sale a Mi: i due suoni formano l'ottava Mi--Mi sulla dominante.
	\musica{es_cap22_03}
\end{esempio}

\begin{esempio}[Sesta tedesca in La minore]
	\label{ex:tedesca_minore}
	La sesta tedesca passa per il 6/4 di cadenza, in cui Fa e Re$\sharp$ sono già arrivati al Mi.
	\musica{es_cap22_04}
\end{esempio}

\begin{esempio}[Seste italiana e francese in Do maggiore]
	\label{ex:italiana_francese_maggiore}
	Le seste eccedenti sono formate sul La$\flat$ (VI grado abbassato): La$\flat$--Fa$\sharp$ si allargano a Sol--Sol.
	\musica{es_cap22_05}
\end{esempio}

\begin{esempio}[Sesta tedesca in Do maggiore]
	\label{ex:tedesca_maggiore}
	Nella sesta tedesca il Mi$\flat$ sale a Mi$\natural$ nel passaggio al 6/4 di tonica maggiore; poi il Mi scende al Re della dominante.
	\musica{es_cap22_06}
\end{esempio}

\begin{errore}[Sesta tedesca che risolve direttamente sulla dominante]
	\label{err:tedesca_diretta}
	Fa $\searrow$ Mi (basso) e Do $\searrow$ Si formano \textbf{quinte parallele} tra tenore e basso (Fa--Do $\rightarrow$ Mi--Si): per evitarle si passa per il 6/4 di cadenza, come negli esempi precedenti.
	\musica{es_cap22_07}
\end{errore}

\section{Quadro riassuntivo}
\label{sec:quadro_seste}

\begin{center}
\small
\begin{tabular}{|l|l|l|p{6cm}|}
\hline
Specie & Note (La minore) & Cifratura & Risoluzione \\
\hline
Italiana & Fa -- La -- Re$\sharp$ (La raddoppiato) & $\sharp$6 & Fa $\searrow$ Mi, Re$\sharp$ $\nearrow$ Mi, La $\searrow$ Sol$\sharp$, La $\nearrow$ Si: V completa \\
Francese & Fa -- La -- Si -- Re$\sharp$ & $\sharp$6 4 3 & Fa $\searrow$ Mi, Re$\sharp$ $\nearrow$ Mi, La $\searrow$ Sol$\sharp$, Si ferma: V completa \\
Tedesca & Fa -- La -- Do -- Re$\sharp$ & $\sharp$6 5 3 & Fa $\searrow$ Mi, Re$\sharp$ $\nearrow$ Mi, La e Do fermi: \sigla{i}{6}{4}; poi La $\searrow$ Sol$\sharp$, Do $\searrow$ Si: V \\
\hline
\end{tabular}
\end{center}

La forma tipica ha il VI$\flat$ al basso. La sesta eccedente appare, nella Regola dell'Ottava, sul VI grado discendente del modo minore (capitolo \ref{chap:regola_ottava_completa}).
